\documentclass[11pt,a4paper]{article}
\usepackage{amsmath,amssymb,amsthm}
\usepackage[margin=2.5cm]{geometry}
\usepackage{hyperref}

\newtheorem{definition}{Definition}[section]
\newtheorem{theorem}[definition]{Theorem}
\newtheorem{proposition}[definition]{Proposition}
\newtheorem{lemma}[definition]{Lemma}
\newtheorem{corollary}[definition]{Corollary}
\newtheorem{conjecture}[definition]{Conjecture}
\newtheorem{remark}[definition]{Remark}

\title{Hyperseed Formalization:\\
  DeAI After Mythos}
\author{ProtomegaTron (automated formalization)}
\date{2026-09-18}

\begin{document}
\maketitle

\begin{abstract}
We formalize the intellectual structure of Ben Goertzel's Substack article
``DeAI After Mythos'' (2026-04-15) within the Hyperseed ontology. The
article analyzes the implications of Anthropic's Mythos cybersecurity
breakthrough for decentralized AI. Key mappings: friction-based security
as base-level obscurity, formal reasoning as fiber-level security,
OmegaClaw as reasoning fiber over discovery base, the post-Mythos regime
as base-level security collapse, decentralized security as distributed
fiber, and the regime change as a fiber transition from base-obscurity
to formal-guarantee security.
\end{abstract}

\tableofcontents

%% =================================================================
\section{Base-level obscurity vs.\ fiber-level security}
\label{sec:base-fiber}
%% =================================================================

\begin{theorem}[The friction model --- claims 1--5, 21, 25]
\label{thm:friction}
Historical software security rested on a friction barrier: severe
vulnerabilities existed in virtually all critical software, but the
effort required to find them created a cost barrier that limited
exploitation. The article's metaphor: a dog guarding its food bowl
--- not robust access control, just that most challengers couldn't
be bothered.

In Hyperseed terms, this is \emph{base-level obscurity}: security
that depends on the difficulty of base-level operations (reading
large codebases, tracking control-flow interactions, correlating
patches with latent vulnerabilities). The security property lives
in the base (discovery difficulty), not in the fiber (structural
guarantee):
\[
  \text{Security}_{\text{friction}} = f(\text{cost of discovery})
  \qquad \text{not} \qquad
  \text{Security}_{\text{formal}} = f(\text{proof of absence})
\]

Base-level obscurity was never real security --- it was an economic
barrier masquerading as a structural property. The ``dog guarding
food bowl'' metaphor captures this precisely: the defense operates
at the base level (physical presence, aggression display, cost of
challenge) and collapses when the economics change.

Mythos changes the economics: frontier models can read very large
codebases, track subtle control-flow interactions, correlate patches
with vulnerabilities, and iterate exploit ideas hundreds of times ---
all at dramatically lower cost than human specialists. The base-level
operations that constituted the friction barrier are now cheap.

The collapse is not gradual but regime-like: once the cost drops
below the threshold where exploitation becomes economically viable
for a wide range of actors, the entire friction-based security model
fails simultaneously across all software that relies on it.
\end{theorem}

%% =================================================================
\section{Fiber-level security through formal reasoning}
\label{sec:formal}
%% =================================================================

\begin{proposition}[Formal methods --- claims 6--10, 22]
\label{prop:formal}
Once raw vulnerability discovery gets cheaper, the scarce resource
becomes trustworthy interpretation and reliable containment: Which
candidate findings are real? Which are exploitable in the actual
deployed environment? Which touch critical assets?

In Hyperseed terms, formal reasoning provides \emph{fiber-level
security}: structural guarantees that hold regardless of how cheap
discovery becomes:
\begin{itemize}
\item \textbf{Base-level security:} ``this bug is hard to find''
  (depends on discovery cost, collapses when cost drops).
\item \textbf{Fiber-level security:} ``this component is formally
  verified to be free of this class of bugs'' (depends on proof
  structure, holds regardless of discovery cost).
\end{itemize}

The shift from base to fiber is the key structural change:
\[
  \text{Value}_{\text{pre-Mythos}} \propto \text{discovery capability}
  \qquad \xrightarrow{\text{Mythos}} \qquad
  \text{Value}_{\text{post-Mythos}} \propto \text{interpretation capability}
\]

Discovery becomes commoditized (many models can do it); interpretation
remains scarce (requires formal reasoning, context understanding,
deployment knowledge). The value chain shifts from base operations
(finding bugs) to fiber operations (understanding, validating,
containing findings).

Mythos doesn't make formal methods obsolete --- it makes them more
necessary. This is a general pattern: when base-level operations
become cheap, fiber-level operations become more valuable. The
same pattern appears in:
\begin{itemize}
\item \textbf{Information:} when information access becomes cheap
  (internet), interpretation and judgment become more valuable.
\item \textbf{Content:} when content creation becomes cheap (LLMs),
  curation and quality assessment become more valuable.
\item \textbf{Security:} when vulnerability discovery becomes cheap
  (Mythos), formal verification and containment become more valuable.
\end{itemize}

In each case, the base operation is commoditized and the fiber
operation captures the value.
\end{proposition}

%% =================================================================
\section{OmegaClaw as reasoning fiber}
\label{sec:omegaclaw}
%% =================================================================

\begin{definition}[Reasoning over discovery --- claims 11--14, 23]
\label{def:omegaclaw}
OmegaClaw is positioned as a reasoning and validation layer: it
provides the interpretation/reasoning fiber over the raw discovery
base. Mythos-class models operate at the base level (finding
vulnerabilities); OmegaClaw operates at the fiber level (interpreting,
validating, containing).

Three fiber-level capabilities:
\begin{enumerate}
\item \textbf{MeTTa reasoning:} formal logical reasoning about software
  --- exactly what becomes scarce and valuable post-Mythos. Not just
  pattern-matching on code but structured logical inference about
  program behavior.
\item \textbf{Hyperon integration:} proto-AGI components provide
  deeper reasoning than pure LLM-based security tools. The reasoning
  is not just statistical (``this pattern looks like a bug'') but
  structural (``this control flow violates this invariant'').
\item \textbf{Decentralized architecture:} distributed security
  reasoning without single-point-of-failure. No single entity controls
  the interpretation layer --- the fiber is distributed.
\end{enumerate}

The architectural insight: security reasoning is a fiber operation,
and fiber operations should be distributed (for robustness) and
formal (for reliability). OmegaClaw combines both: distributed
architecture with formal reasoning capability.

This positions OmegaClaw as the fiber layer in a two-layer security
architecture:
\[
  \underbrace{\text{Mythos-class models}}_{\text{Base: discovery}}
  \quad + \quad
  \underbrace{\text{OmegaClaw}}_{\text{Fiber: interpretation}}
  \quad = \quad
  \text{Full security stack}
\]
\end{definition}

%% =================================================================
\section{The regime change as fiber transition}
\label{sec:regime}
%% =================================================================

\begin{theorem}[Fiber transition --- claims 18--20, 24, 26]
\label{thm:regime}
The Mythos regime change is a fiber transition: the security fiber
of the software ecosystem is shifting from base-obscurity fiber to
formal-guarantee fiber.

Three phases:
\begin{enumerate}
\item \textbf{Pre-Mythos:} base-obscurity fiber dominant. Security
  rests on friction (bugs hard to find). Formal methods are nice-to-have,
  not essential.
\item \textbf{Transition (now):} base-obscurity collapsing. Systems
  still relying on friction are catastrophically exposed. Formal
  methods become essential but are not yet widely deployed. This is
  the crisis period.
\item \textbf{Post-Mythos:} formal-guarantee fiber dominant. Security
  rests on structural guarantees (formal verification, proof-carrying
  code). Discovery is commoditized; interpretation and containment
  are the value center.
\end{enumerate}

The transition period is where crisis and opportunity coexist:
\begin{itemize}
\item \textbf{Crisis:} more vulnerabilities discoverable by more actors
  on faster timescales. Every system relying on base-obscurity is
  exposed.
\item \textbf{Opportunity:} reasoning systems that can interpret,
  validate, and contain findings become essential infrastructure.
  Projects with formal reasoning capabilities capture the value
  that shifts from base to fiber.
\end{itemize}

This is a security analog of the scalar collapse principle: base-level
security was always an illusion --- it looked like security but was
really just cost economics. Only fiber-level security (structural
guarantees) was ever real. Mythos didn't create new insecurity; it
revealed the insecurity that was always there by removing the cost
barrier that obscured it.

For ASI:Chain specifically, the post-Mythos bar means:
\begin{itemize}
\item Formal verification of critical components (not optional).
\item Narrower, well-defined attack surfaces (not sprawling interfaces).
\item Explicit trust boundaries (not implicit assumptions).
\item Vague ``we take security seriously'' is insufficient when
  frontier models can probe every surface.
\end{itemize}
\end{theorem}

%% =================================================================
\section{Distributed security fiber}
\label{sec:distributed}
%% =================================================================

\begin{proposition}[Decentralized security --- claim 27]
\label{prop:distributed}
OmegaClaw's decentralized architecture provides distributed security
fiber: no single point of failure, no single point of capture, no
single entity controlling the interpretation layer.

In Hyperseed terms, this is the security application of the
diversity-as-safety principle:
\begin{itemize}
\item \textbf{Centralized security reasoning:} single fiber, single
  failure mode, single capture point. If the centralized security
  service is compromised, all dependent systems lose their security
  fiber simultaneously.
\item \textbf{Distributed security reasoning:} diverse fiber bundle,
  uncorrelated failure modes, no single capture point. Individual
  nodes can be compromised without collapsing the entire security
  fiber.
\end{itemize}

This is especially important post-Mythos: if frontier models can
probe every surface, a centralized security service is itself a
high-value target. A distributed security fiber is harder to
compromise comprehensively --- the attacker must compromise many
independent nodes rather than one central service.
\end{proposition}

%% =================================================================
\section{Cross-references}
%% =================================================================

\begin{remark}[Connections to other formalizations]
\begin{itemize}
\item \textbf{Bootleggers \& Baptists} (2026-04-28): regulatory
  capture. The centralized security model is vulnerable to the same
  capture dynamics as centralized regulation --- whoever controls
  the security service controls the interpretation of vulnerabilities.
\item \textbf{OpenBGI} (2026-05-07): decentralized infrastructure.
  OmegaClaw's decentralized security architecture parallels OpenBGI's
  decentralized governance.
\item \textbf{Seeding RSI} (2026-07-28): Hyperon architecture.
  OmegaClaw's MeTTa reasoning and Hyperon integration connect
  directly to the Hyperon proto-AGI platform.
\item \textbf{The OpenAI/Microsoft Shakeup} (2026-05-01): structural
  independence. OmegaClaw's decentralized architecture provides
  the structural independence from any single provider that OpenAI
  failed to maintain.
\item \textbf{Playing Around with ARC-AGI-3} (2026-05-08): verification.
  The LLM+verifier pattern (generate hypotheses, verify formally)
  parallels the Mythos+OmegaClaw pattern (discover vulnerabilities,
  verify formally).
\end{itemize}
\end{remark}

\begin{conjecture}[Base-to-fiber value migration]
Conjecture: every domain undergoes a base-to-fiber value migration
when the base operations become commoditized by AI. The pattern:
\begin{enumerate}
\item Pre-AI: base operations are expensive and scarce. Value lives
  in performing the base operations (finding bugs, creating content,
  generating code).
\item AI commoditization: base operations become cheap. Value shifts
  to fiber operations (interpreting, validating, containing, curating).
\item Post-migration: fiber operations dominate value. Base operations
  are infrastructure, not value centers.
\end{enumerate}

If this conjecture holds, the Mythos example is just the security
instance of a universal pattern. Every domain where AI commoditizes
base operations will see the same value migration to fiber operations.
The organizations and systems that capture the fiber layer will
capture the value; those that remain at the base layer will be
commoditized.
\end{conjecture}

\end{document}
